\section{Elementary arithmetic correspondences}\label{sec:models}
The models in Theorem~\ref{thm:arithmetic-exhaustion} have an explicit
counting-space realization. We give the scalar transport and the
normality proof before beginning the geometric argument; neither depends
on arithmetic exhaustion. Projective induction is classical
\cite[Section~4]{Mackey1958}, and these coordinates also belong to the
finite-index bimodule calculus of \cite[Section~6.2]{Vaes2008}.

Throughout this section $\Gamma,\Lambda$ are arbitrary countably infinite
ICC groups, $M=L_\mu(\Gamma)$, and $N=L_\omega(\Lambda)$. Put
\[
 D=\Gamma\times\Lambda,\qquad
 \alpha((g,h),(g',h'))=\mu(g,g')\overline{\omega(h,h')}.
\]
For an isomorphism $\delta:A\to B$ of finite-index subgroups, its graph is
$C_\delta=\{(a,\delta(a)):a\in A\}\leq D$.

\subsection{Projective induction}
Let $C\leq D$ be a subgroup and let $\gamma$ be a normalized unitary
projective representation on a finite-dimensional space $V$, with
multiplier $\alpha|_{C\times C}$.  Write $X=D/C$ for left cosets.
Choose representatives $r_x\in D$, with $r_C=e$, and give each $x\in X$
a copy $V_x$ of $V$.  Set
$\mathcal H(C,\gamma)=\bigoplus_{x\in X}V_x$, with the counting
Hilbert norm.  For $s\in D$ put
\begin{equation}\label{eq:input-A1}
 \begin{split}
 c(s,x)&=r_{sx}^{-1}sr_x,\\
 U^X_s(x)&=
 \frac{\alpha(s,r_x)}{\alpha(r_{sx},c(s,x))}\,
 \gamma(c(s,x)):V_x\longrightarrow V_{sx}.
 \end{split}
\end{equation}

\begin{lemma}[Exact projective induction]
\label{lem:models-induction}
Equation~\eqref{eq:input-A1} defines a unitary projective representation
of $D$ with multiplier $\alpha$.  Changing the coset representatives
gives a unitarily equivalent bundle by unitary identifications of its
individual fibres.  If $C$ is the graph of a finite-index subgroup
isomorphism $A\to B$, the bundle has commuting normal $M$ and right
$N$ actions, and is bifinite.
\end{lemma}

\begin{proof}
The correction $c(s,x)$ belongs to $C$, and
\[
 c(s,tx)c(t,x)=c(st,x).
\]
Using the cocycle identity for $\alpha$ and the projective multiplication
of $\gamma$ in Equation~\eqref{eq:input-A1} gives
\begin{equation}\label{eq:calculus-transport-identity}
 U^X_s(tx)U^X_t(x)=\alpha(s,t)U^X_{st}(x).
\end{equation}
One can also check all the scalar factors at once in the central
extension with multiplication
$(z,s)(w,t)=(zw\alpha(s,t),st)$: the representation of its preimage
of $C$ is $(z,c)\mapsto z\gamma(c)$, and the equality
$sr_x=r_{sx}c(s,x)$ gives exactly the ratio in
Equation~\eqref{eq:input-A1}.  Thus this argument retains, rather than
discards, the fixed multiplier.

For explicit independence of representatives, suppose
$r'_x=r_xk_x$, where $k_x\in C$.  The identification from the new
coordinate copy of $V_x$ to the old one is
\[
 F_x=\alpha(r_x,k_x)^{-1}\gamma(k_x).
\]
Indeed identifying a fibre by $U_{r'_x}$ instead of $U_{r_x}$ has
precisely this effect.  Equation~\eqref{eq:calculus-transport-identity}
then gives $F_{sx}U'{}^X_s(x)=U^X_s(x)F_x$.

Now assume that $C$ is a graph.  Both coordinate groups act freely on
$X$.  For example, if $(g,e)r_xC=r_xC$, then
$r_x^{-1}(g,e)r_x\in C$ has second coordinate $e$, so $g=e$.
On an orbit of the first coordinate group choose a fibre $V_{x_0}$.
The unitary identification
\[
 \ell^2(\Gamma)\otimes V_{x_0}\longrightarrow
 \bigoplus_{g\in\Gamma}V_{(g,e)x_0},
 \qquad \delta_g\otimes\xi\longmapsto U^X_{(g,e)}\xi
\]
intertwines its action with $\lambda_\mu\otimes1$, since
$U_{(a,e)}U_{(g,e)}=\mu(a,g)U_{(ag,e)}$.  Hence it extends normally
to $L_\mu(\Gamma)$.  The same argument for the second coordinate
gives copies of $\lambda_{\bar\omega}$, which are the standard
normal representations of the generators acting by right adjoints
$\xi\mapsto\xi v_h^*$.  The two actions commute, because the product
cocycle has trivial mixed-coordinate factors.

The first coordinate orbits are indexed by $\Lambda/B$, and the
second coordinate orbits by $\Gamma/A$.  They are finite in number.
Each is a regular module with multiplicity $\dim_{\mathbb C}V$,
proving bifiniteness as well as the dimension formulas below.
\end{proof}

We next identify this construction with the Connes-product model in
Theorem~\ref{thm:arithmetic-exhaustion}.  This also supplies a direct normality
check of the subgroup action used to define that product.

\begin{lemma}[The graph-bundle unitary]
\label{lem:models-graph}
Let $\delta:A\to B$ be a finite-index subgroup isomorphism and let
$\sigma:A\to\mathcal U(V)$ have multiplier
$\beta=\mu|_{A\times A}/\delta^*\omega$.  Define
$\gamma(a,\delta(a))=\sigma(a)$.  Then
\[
 E(A,B,\delta,\sigma)=
 L^2(M)\mathbin{\bar\otimes}_{L_{\mu|A}(A)}
       (V\otimes L^2(N))
 \cong\mathcal H(C_\delta,\gamma).
\]
The unitary sends the vector $\xi$ in the fibre labelled by $x$ to
$U_{r_x}j\xi$, where
$j\xi=\widehat1\otimes(\xi\otimes\widehat1)$ in the displayed model.
\end{lemma}

\begin{proof}
On $V\otimes L^2(N)$ the subgroup action is
\[
 u_a(\xi\otimes\eta)=\sigma(a)\xi\otimes v_{\delta(a)}\eta.
\]
Its multiplier is $\beta(a,b)\omega(\delta(a),\delta(b))=\mu(a,b)$.
For normality, choose representatives $t$ for $B\backslash\Lambda$.
On the summand spanned by $\xi\otimes v_{\delta(a)}v_t$, the map
\begin{equation}\label{eq:calculus-absorption}
 \xi\otimes v_{\delta(a)}v_t
 \longmapsto\delta_a\otimes\sigma(a)^*\xi
 \quad\text{in }\ell^2(A)\otimes V
\end{equation}
is unitary and conjugates the subgroup action to
$\lambda_{\mu|A}\otimes1$.  To check the phase, applying $u_b$
before this map produces
\[
 \omega(\delta(b),\delta(a))
 \sigma(ba)^*\sigma(b)
 =\omega(\delta(b),\delta(a))\beta(b,a)\sigma(a)^*
 =\mu(b,a)\sigma(a)^*.
\]
Thus the left subgroup action is normal.  Its right $N$ action is
the standard action on the second factor and commutes with it.

Choose representatives $r$ for $\Gamma/A$.  As a right
$L_{\mu|A}(A)$ module, $L^2(M)$ has orthonormal free basis $u_r$;
the assertion follows directly from
$E_{L_{\mu|A}(A)}(u_r^*u_{r'})=\delta_{r,r'}1$.
The relative product therefore identifies with the orthogonal sum
\begin{equation}\label{eq:calculus-coset-model}
 E(A,B,\delta,\sigma)
 =\bigoplus_{r\in\Gamma/A}u_r\otimes(V\otimes L^2(N)).
\end{equation}
At its identity fibre the graph action is exactly
\[
 U_{(a,\delta(a))}j\xi
 =u_a\otimes(\xi\otimes v_{\delta(a)}^*)
 =j\sigma(a)\xi.
\]
For $g=ra$, balancing in Equation~\eqref{eq:calculus-coset-model}
gives the more explicit formula
\[
 U_{(g,h)}j\xi
 =\mu(r,a)^{-1}u_r\otimes
       \bigl(\sigma(a)\xi\otimes v_{\delta(a)}v_h^*\bigr).
\]
The last Fourier coordinate is $\delta(a)h^{-1}$, up to its displayed
unit scalar.  Thus two translated identity fibres are orthogonal
unless their labels in $D/C_\delta$ agree: equality of their first
cosets and last Fourier coordinates is exactly that coset equality.
Conversely these fibres exhaust every summand in
Equation~\eqref{eq:calculus-coset-model}.  They are an orthogonal,
transitive system with graph action $\gamma$ at the basepoint.
Finally
\[
 U_sU_{r_x}j
 =\frac{\alpha(s,r_x)}{\alpha(r_{sx},c(s,x))}
       U_{r_{sx}}j\gamma(c(s,x)),
\]
which proves that the stated unitary has exactly the transport
of Equation~\eqref{eq:input-A1}.
\end{proof}


\begin{corollary}[Elementary normality and dimensions]\label{cor:models-dimensions}
The correspondence $E(A,B,\delta,\sigma)$ is normal and bifinite. If
$n=\dim_{\mathbb C}V_\sigma$, then
\[
 \dim_M E(A,B,\delta,\sigma)=n[\Lambda:B],\qquad
 \dim E(A,B,\delta,\sigma)_N=n[\Gamma:A].
\]
\end{corollary}
\begin{proof}
Use Lemma~\ref{lem:models-graph} and the orbit decompositions in
Lemma~\ref{lem:models-induction}. Each free coordinate orbit gives $n$
regular copies, of module dimension $n$ with the unnormalized matrix
trace convention. The respective orbit counts are $[\Lambda:B]$ and
$[\Gamma:A]$.
\end{proof}

Finite direct sums and closed bimodule summands of these models remain
normal and bifinite: the two actions restrict normally to a reducing
subspace, and module dimension is additive and monotone under orthogonal
projections in the commutant
\cite[Theorem~10.2.1]{Jones2009}. The geometric problem is the converse:
starting from an arbitrary bifinite correspondence, recover enough
subgroup graphs to contain it. We now make the finite-index and compact
reductions needed for that construction.
